% ===========================================================================
%  sviluppatore/08-qualita/01-prove.tex
% ===========================================================================
\section{Come sono scritte le prove}\label{sez:prove}

Non c'è un \pkg{pytest} con mille asserzioni minute. Ci sono cinque suite,
ciascuna dedicata a un \emph{contratto} del sistema, che si lanciano come
script e stampano un resoconto leggibile. La scelta ha una conseguenza utile:
ogni suite si apre con un docstring che dice quale problema sorveglia, e
quello è il posto giusto dove cercare prima di modificare l'area che copre.

\begin{center}
\begin{tabular}{@{}L{48mm}R{14mm}L{76mm}@{}}
\toprule
\textbf{Suite} & \textbf{Righe} & \textbf{Che contratto sorveglia} \\
\midrule
\rigaTenue
\file{prova\_normalizzazione\_stream.py} & \num{691}
  & Backtest e live devono produrre lo \emph{stesso libro}
    (\vediSez{live}). \\
\file{prova\_beta\_accesso.py} & \num{614}
  & L'accesso riservato in beta, \textbf{contro il sito vero}. \\
\rigaTenue
\file{prova\_covarianza.py} & \num{489}
  & Lo shrinkage di Ledoit--Wolf e la modalità \cod{covariance}. \\
\file{prova\_pesi\_dichiarati.py} & \num{484}
  & Il contratto di \cod{Signal.target\_weight} e la modalità
    \cod{explicit}. \\
\rigaTenue
\file{prova\_permessi\_marketplace.py} & \num{429}
  & I due strati del permesso di pubblicare, \textbf{contro il sito vero}. \\
\bottomrule
\end{tabular}
\end{center}

\subsection{Due suite girano contro il sito vero}

Non usano finti. Avviano un Next.js e un Postgres reali e interrogano gli
endpoint veri. È più lento e più fastidioso da far partire, e la ragione è
dichiarata nel docstring: i permessi e l'accesso non vivono nel codice
Python, vivono nelle regole del sito. Una prova con dei finti verificherebbe
che la dashboard sa leggere una risposta --- cosa che non è mai stata il
problema --- e non che la risposta è quella giusta.

\subsection{Le proprietà che una prova deve avere qui}

Tre, e tutte e tre sono state imparate rompendole.

\begin{description}[leftmargin = 1.6em, style = nextline, itemsep = 6pt]

\item[Ripetibile]
  Una suite che modifica lo stato (prezzi dei piani, elenchi di ammessi) deve
  azzerarlo all'inizio, non alla fine: se una esecuzione si interrompe a metà,
  la pulizia finale non viene eseguita e la successiva parte da uno stato
  sporco. \file{prova\_permessi\_marketplace.py} comincia con una
  reimpostazione in SQL; \file{prova\_beta\_accesso.py} ha una funzione
  \simbolo{azzera\_prova()} e un elenco fisso di indirizzi.

\item[Senza soglie indovinate]
  Un'asserzione come \cod{assert scala < 1.3} è una soglia inventata: passa o
  fallisce per ragioni che nessuno ha verificato. La forma corretta è
  relativa --- \cod{assert scala\_degenere < 0.75 * scala\_corretta} --- perché
  confronta due cose che il sistema produce, invece di confrontarne una con
  un numero scritto a mano.

\item[Che isola ciò che misura]
  Misurare la proporzionalità confrontando due chiamate complete non
  funziona se ciascuna chiamata registra un'osservazione in più e cambia lo
  stato. Il confronto va fatto sulla funzione che si vuole misurare, da sola.

\end{description}

\begin{attenzione}{Una prova che non è ripetibile è peggio di nessuna prova}
Passa la prima volta, fallisce la seconda, e chi la vede fallire conclude
che la prova è inaffidabile --- non che il codice lo sia. Da quel momento il
risultato di quella suite viene ignorato, cioè la suite ha smesso di
esistere pur continuando a girare.
\end{attenzione}

\subsection{Il resto di \file{tests/}}

Oltre alle cinque suite ci sono: prove end-to-end (\file{e2e\_hedging.py},
\file{e2e\_regime\_modulation.py}), confronti A/B fra strategie, diagnosi
specifiche per singole strategie (\file{diagnose\_ising.py},
\file{diagnose\_regime\_signals.py}), un server SMTP finto
(\file{posta\_finta.py}) usato dalla prova dell'accesso, e verifiche
numeriche di migrazione (\file{verify\_migration\_numbers.py}).
